\begin{figure}[htbp]
    \centering
    \begin{tikzpicture}[
        block/.style={
            draw,
            rectangle,
            rounded corners=3pt,
            minimum width=2.25cm,
            minimum height=1.15cm,
            align=center,
            fill=white,
            line width=0.8pt,
            font=\scriptsize\sffamily\bfseries
        },
        arrow/.style={-{Stealth[scale=1.15]}, line width=0.8pt},
        note/.style={font=\tiny\sffamily, align=center}
    ]
        \node (csr) [block] {Cochlear\\Spike Raster};
        \node (feature) [block, right=0.42cm of csr] {Feature\\Extraction};
        \node (compare) [block, right=0.42cm of feature] {Cue\\Comparison};
        \node (population) [block, right=0.42cm of compare] {Candidate\\Population};
        \node (readout) [block, right=0.42cm of population] {SC-Style\\Readout};

        \draw [arrow] (csr) -- (feature);
        \draw [arrow] (feature) -- (compare);
        \draw [arrow] (compare) -- (population);
        \draw [arrow] (population) -- (readout);

        \node [note, below=0.22cm of feature] {onset timing or\\spectral profile};
        \node [note, below=0.22cm of compare] {corollary discharge,\\opposite ear, or template};
        \node [note, below=0.22cm of population] {activity over candidate\\coordinate values};
    \end{tikzpicture}
    \caption{Shared computational pattern used by the three localisation pathways. Each pathway preserves an interpretable intermediate population rather than collapsing the received waveform directly into a scalar prediction.}
    \label{fig:common_pathway_pattern}
\end{figure}
